\documentclass[11pt,a4paper]{article}

\usepackage[a4paper,margin=2.5cm]{geometry}
\usepackage{amsmath,amssymb,bm}
\usepackage{physics}
\usepackage{siunitx}
\usepackage{hyperref}
\usepackage{enumitem}

\hypersetup{
  colorlinks=true,
  linkcolor=blue,
  citecolor=blue,
  urlcolor=blue
}

\newcommand{\ii}{\mathrm{i}}
\newcommand{\dd}{\mathrm{d}}
\newcommand{\kvec}{\mathbf{k}}
\newcommand{\qvec}{\mathbf{q}}
\newcommand{\rvec}{\mathbf{r}}
\newcommand{\Rvec}{\mathbf{R}}
\newcommand{\epsvec}{\boldsymbol{\epsilon}}
\newcommand{\Ylm}{Y_{lm}}
\newcommand{\mt}{\mathrm{MT}}
\newcommand{\Ha}{\mathrm{Ha}}
\newcommand{\eV}{\mathrm{eV}}
\newcommand{\absr}{a}
\newcommand{\sig}{\sigma}
\newcommand{\tausp}{\tau}

\title{Independent-Particle X-ray Absorption Spectroscopy in FLEUR:\\
Implemented Formalism, Notation, and Limitations}
\author{Implementation notes for the FLEUR XAS prototype}
\date{\today}

\begin{document}

\maketitle

\begin{abstract}
These notes summarize the independent-particle X-ray absorption spectroscopy
(XAS) formalism implemented in the FLEUR prototype. The implementation evaluates
electric-dipole transitions from relativistic atomic core states localized on a
chosen absorber atom into unoccupied Kohn--Sham or generalized Kohn--Sham band
states. The current implementation supports the separate \(K\), \(L_2\), and
\(L_3\) edges, linear polarizations \(x,y,z\), scalar calculations, and the
validated first-variation noncollinear spin-orbit-coupled path. The purpose of
this document is twofold: first, to define the implemented equations and all
indices explicitly; second, to state clearly what physical effects are and are
not included in the present independent-particle approximation.
\end{abstract}

\section{Physical quantity}

X-ray absorption spectroscopy measures the probability that an incoming photon
promotes a deep core electron into an available unoccupied state. In the present
implementation, the photon momentum is neglected in the transition operator
and the electric-dipole approximation is used. The photon is therefore described
only by its polarization vector
\begin{equation}
  \epsvec = (\epsilon_x,\epsilon_y,\epsilon_z),
\end{equation}
and the dipole operator is
\begin{equation}
  \hat{D}_{\epsvec} = \epsvec \cdot \rvec .
\end{equation}

The implemented spectrum for a polarization \(\epsvec\) is
\begin{equation}
\label{eq:xas-spectrum}
  I_{\epsvec}(\omega)
  =
  \sum_{\absr\in A_Z}
  \sum_{c\in E}
  \sum_{m_j=-j_c}^{j_c}
  \sum_{n,\kvec}
  w_{\kvec}\,
  \bigl[1-f_{n\kvec}\bigr]\,
  \left|
     M^{\absr}_{n\kvec,c m_j}(\epsvec)
  \right|^2
  g_{\eta}\!\left[
     \omega -
     \left(\varepsilon_{n\kvec}-\varepsilon^{\absr}_{c}\right)
  \right].
\end{equation}
All indices in Eq.~\eqref{eq:xas-spectrum} are as follows.
\begin{description}[leftmargin=2.7cm,style=nextline]
  \item[\(\absr\)] absorber atom. The implementation selects all atoms with
  the requested nuclear charge \(Z\), denoted collectively by \(A_Z\).
  \item[\(A_Z\)] set of absorber atoms with nuclear charge \(Z\). For example,
  for \texttt{absorberZ="26"}, \(A_Z\) is the set of Fe atoms.
  \item[\(E\)] selected absorption edge. The currently implemented separate
  edges are \(K\), \(L_2\), and \(L_3\).
  \item[\(c\)] core level associated with the selected edge. For a fixed edge,
  \(c\) specifies the principal quantum number \(n_c\), orbital angular momentum
  \(l_c\), total angular momentum \(j_c\), and relativistic quantum number
  \(\kappa_c\).
  \item[\(m_j\)] magnetic quantum number of the core state, with
  \(m_j=-j_c,-j_c+1,\ldots,j_c\).
  \item[\(n\)] band index of the final Kohn--Sham or generalized Kohn--Sham
  state.
  \item[\(\kvec\)] Bloch wave vector in the Brillouin zone.
  \item[\(w_{\kvec}\)] integration weight of the \(k\)-point.
  \item[\(f_{n\kvec}\)] occupation of band \(n\) at \(k\)-point \(\kvec\).
  The factor \(1-f_{n\kvec}\) restricts the spectrum to available final states.
  \item[\(\varepsilon_{n\kvec}\)] band eigenvalue of the final state.
  \item[\(\varepsilon^{\absr}_{c}\)] core eigenvalue for core level \(c\) on
  absorber atom \(\absr\).
  \item[\(M^{\absr}_{n\kvec,c m_j}(\epsvec)\)] dipole transition matrix element.
  \item[\(g_{\eta}\)] normalized broadening function with width parameter
  \(\eta\). In the current code this is the numerical broadening used to turn
  the discrete transition list into a smooth spectrum.
  \item[\(\omega\)] spectrum energy variable. In the code output this is the
  transition energy on the chosen energy grid.
\end{description}

Equation~\eqref{eq:xas-spectrum} is an independent-particle form of Fermi's
golden rule. It uses the electronic structure from the ground-state calculation
(or from whatever self-consistent state the user has supplied), and it does not
solve a new interacting core-excited many-body problem.

\section{Supported edges and core quantum numbers}

The current implementation treats \(K\), \(L_2\), and \(L_3\) as separate
edges. A combined \(L_{23}\) mode is deliberately not implemented yet; users
should run \(L_2\) and \(L_3\) separately.

The supported edge mappings are
\begin{center}
\begin{tabular}{c c c c c c}
\hline
Edge & Core shell & \(n_c\) & \(l_c\) & \(j_c\) & \(\kappa_c\) \\
\hline
\(K\)   & \(1s_{1/2}\) & 1 & 0 & \(1/2\) & \(-1\) \\
\(L_2\) & \(2p_{1/2}\) & 2 & 1 & \(1/2\) & \(+1\) \\
\(L_3\) & \(2p_{3/2}\) & 2 & 1 & \(3/2\) & \(-2\) \\
\hline
\end{tabular}
\end{center}
The relativistic quantum number is defined by the usual Dirac convention:
\begin{equation}
  \kappa =
  \begin{cases}
    -(l+1), & j=l+\frac12,\\[3pt]
    +l, & j=l-\frac12 .
  \end{cases}
\end{equation}

\section{Dipole matrix element}

The absorber-resolved dipole matrix element in Eq.~\eqref{eq:xas-spectrum} is
\begin{equation}
\label{eq:matrix-element}
  M^{\absr}_{n\kvec,c m_j}(\epsvec)
  =
  \left\langle
     \Psi_{n\kvec}
  \middle|
     \epsvec\cdot\rvec
  \middle|
     \Phi^{\absr}_{c m_j}
  \right\rangle_{\mt_{\absr}} .
\end{equation}
The integral is evaluated inside the muffin-tin sphere \(\mt_{\absr}\) around
absorber atom \(\absr\). The symbols are:
\begin{description}[leftmargin=2.7cm,style=nextline]
  \item[\(\Psi_{n\kvec}\)] final band state.
  \item[\(\Phi^{\absr}_{c m_j}\)] core state localized on absorber atom \(\absr\).
  \item[\(\mt_{\absr}\)] muffin-tin sphere of absorber atom \(\absr\).
\end{description}

In the muffin-tin basis the final state is expanded in angular and radial basis
functions. In a spinor calculation we write schematically
\begin{equation}
\label{eq:mt-expansion}
  \Psi_{n\kvec}(\rvec)
  =
  \sum_{\sig}
  \sum_{l m}
  \sum_{\nu}
  A^{\absr,\sig}_{n\kvec,lm\nu}\,
  u^{\absr,\sig}_{l\nu}(r)\,
  Y_{lm}(\hat{\rvec})\,
  \chi_{\sig}.
\end{equation}
Here:
\begin{description}[leftmargin=2.7cm,style=nextline]
  \item[\(\sig\)] spin component. In first-variation noncollinear calculations
  this is a local spin-frame component on the atom.
  \item[\(l,m\)] orbital angular momentum and magnetic quantum number of the
  final-state muffin-tin expansion.
  \item[\(\nu\)] radial basis-function index. It labels the LAPW radial
  solution \(u_l\), its energy derivative \(\dot{u}_l\), and local orbitals.
  \item[\(A^{\absr,\sig}_{n\kvec,lm\nu}\)] muffin-tin expansion coefficient.
  In the FLEUR code these are the corresponding \texttt{abc} coefficients.
  \item[\(u^{\absr,\sig}_{l\nu}(r)\)] radial basis function for angular momentum
  \(l\), radial order \(\nu\), and spin channel \(\sig\).
  \item[\(Y_{lm}\)] complex spherical harmonic.
  \item[\(\chi_{\sig}\)] spin basis function.
\end{description}

The relativistic core state is written as a spin-angular coupled state,
\begin{equation}
\label{eq:core-state}
  \Phi^{\absr}_{c m_j}(\rvec)
  =
  P^{\absr}_{c}(r)
  \sum_{m_c,\sig_c}
  \left\langle
    l_c m_c,\frac12 \sig_c
  \middle|
    j_c m_j
  \right\rangle
  Y_{l_c m_c}(\hat{\rvec})\,\chi_{\sig_c}.
\end{equation}
Here:
\begin{description}[leftmargin=2.7cm,style=nextline]
  \item[\(P^{\absr}_{c}(r)\)] reduced radial core function for core level \(c\).
  \item[\(l_c,m_c\)] core orbital angular momentum and its magnetic quantum
  number.
  \item[\(\sig_c=\pm 1/2\)] core spin projection.
  \item[\(\langle l_c m_c,\frac12\sig_c|j_c m_j\rangle\)] Clebsch--Gordan
  coefficient coupling orbital and spin angular momenta to total angular
  momentum \(j_c,m_j\).
\end{description}

Using Eqs.~\eqref{eq:mt-expansion} and \eqref{eq:core-state}, the matrix element
can be written as
\begin{equation}
\label{eq:M-expanded}
  M^{\absr}_{n\kvec,c m_j}(\epsvec)
  =
  \sum_{\sig}
  \sum_{l m}
  \sum_{\nu}
  \left(A^{\absr,\sig}_{n\kvec,lm\nu}\right)^{*}
  R^{\absr,\sig}_{l\nu,c}
  C^{\sig}_{lm,c m_j}(\epsvec),
\end{equation}
where the three factors are:
\begin{enumerate}[label=(\roman*)]
  \item the complex conjugated band-dependent muffin-tin coefficient
  \(\left(A^{\absr,\sig}_{n\kvec,lm\nu}\right)^{*}\);
  \item the radial integral \(R^{\absr,\sig}_{l\nu,c}\);
  \item the angular dipole coefficient \(C^{\sig}_{lm,c m_j}(\epsvec)\).
\end{enumerate}

\section{Radial integrals}

The radial factor in Eq.~\eqref{eq:M-expanded} is
\begin{equation}
\label{eq:radial-integral}
  R^{\absr,\sig}_{l\nu,c}
  =
  \int_0^{R_{\mt}^{\absr}}
  \dd r\,
  r\,
  U^{\absr,\sig}_{l\nu}(r)\,
  P^{\absr}_{c}(r),
\end{equation}
where
\begin{description}[leftmargin=2.7cm,style=nextline]
  \item[\(R_{\mt}^{\absr}\)] muffin-tin radius of absorber atom \(\absr\).
  \item[\(U^{\absr,\sig}_{l\nu}(r)\)] reduced radial final-state basis function.
  \item[\(P^{\absr}_{c}(r)\)] reduced radial core function.
\end{description}
The use of reduced radial functions is important. With the FLEUR reduced radial
functions, the dipole radial factor is of the form shown in
Eq.~\eqref{eq:radial-integral}.

\section{Angular dipole coefficient}

The Cartesian polarization is converted to spherical components according to
\begin{align}
  \epsilon_{+1} &= -\frac{\epsilon_x+\ii\epsilon_y}{\sqrt{2}},\\
  \epsilon_{0}  &= \epsilon_z,\\
  \epsilon_{-1}&= \frac{\epsilon_x-\ii\epsilon_y}{\sqrt{2}} .
\end{align}
The angular coefficient in Eq.~\eqref{eq:M-expanded} is
\begin{equation}
\label{eq:angular-coefficient}
  C^{\sig}_{lm,c m_j}(\epsvec)
  =
  \sum_{q=-1}^{1}
  \sum_{m_c=-l_c}^{l_c}
  \epsilon_q\,
  G(lm,1q,l_c m_c)\,
  \left\langle
    l_c m_c,\frac12 \sig
  \middle|
    j_c m_j
  \right\rangle .
\end{equation}
Here:
\begin{description}[leftmargin=2.7cm,style=nextline]
  \item[\(q=-1,0,+1\)] spherical component of the dipole operator.
  \item[\(\epsilon_q\)] spherical polarization component.
  \item[\(G(lm,1q,l_c m_c)\)] Gaunt integral,
  \begin{equation}
    G(lm,1q,l_c m_c)
    =
    \int \dd\Omega\,
    Y^*_{lm}(\hat{\rvec})\,
    Y_{1q}(\hat{\rvec})\,
    Y_{l_c m_c}(\hat{\rvec}) .
  \end{equation}
\end{description}
The angular factor enforces the electric-dipole selection rule
\begin{equation}
  \Delta l = l-l_c = \pm 1 .
\end{equation}
Therefore:
\begin{align}
  K\ \mathrm{edge}:  \quad & l_c=0 \rightarrow l=1,\\
  L_2,L_3\ \mathrm{edges}: \quad & l_c=1 \rightarrow l=0,2 .
\end{align}

\section{Noncollinear magnetism and spin-orbit coupling}

In first-variation noncollinear spin-orbit calculations, the band states are
spinors and each atom type can have its own local spin frame. Let
\(U_{\absr}\) be the \(2\times 2\) unitary matrix that maps local spin
components on atom \(\absr\) to global spin components. The FLEUR
\texttt{abc} coefficients in a noncollinear calculation are in the local spin
frame:
\begin{equation}
\label{eq:local-spin}
  \mathbf{A}^{\mathrm{local}}_{\absr}
  =
  U_{\absr}^{\dagger}
  \mathbf{A}^{\mathrm{global}}_{\absr}.
\end{equation}
Here \(\mathbf{A}\) is a two-component spinor of muffin-tin expansion
coefficients for fixed \(n,\kvec,l,m,\nu\). The angular core coefficient must be
expressed in the same local spin frame before being contracted with
\(\mathbf{A}^{\mathrm{local}}_{\absr}\). This local-frame consistency is
essential; otherwise the spectrum can acquire a spurious dependence on the
chosen spin quantization axis.

For symmetry-reduced calculations, FLEUR uses symmetry-generated star members
to reconstruct the full Brillouin-zone contribution. For spinors the
star reconstruction must transform both the orbital part and the spinor part.
For a unitary spatial operation \(R\), let \(S_R\) be the corresponding
spin-\(\frac12\) SU(2) matrix in the global spin frame. The validated
coefficient-space spinor transform is
\begin{equation}
\label{eq:spinor-star}
  \mathbf{A}^{\mathrm{local}}_{\mathrm{star}}
  =
  U_{\mathrm{mapped}}^{\dagger}\,
  S_R^{\dagger}\,
  U_{\mathrm{parent}}\,
  \mathbf{A}^{\mathrm{local}}_{\mathrm{parent}} .
\end{equation}
The indices are:
\begin{description}[leftmargin=2.7cm,style=nextline]
  \item[\(\mathbf{A}^{\mathrm{local}}_{\mathrm{parent}}\)] local spinor
  coefficient at the irreducible parent \(k\)-point.
  \item[\(\mathbf{A}^{\mathrm{local}}_{\mathrm{star}}\)] local spinor
  coefficient reconstructed for the symmetry-generated full-zone star member.
  \item[\(U_{\mathrm{parent}}\)] local-to-global spin-frame matrix for the
  parent absorber atom.
  \item[\(U_{\mathrm{mapped}}\)] local-to-global spin-frame matrix for the
  absorber atom after symmetry mapping.
  \item[\(S_R\)] SU(2) matrix corresponding to the spatial rotation \(R\).
\end{description}
The adjoint \(S_R^\dagger\), rather than \(S_R\), appears because the code
transforms expansion coefficients. Coefficients transform contragrediently to
basis functions. This convention was validated by comparing symmetry-reduced
spectra with explicit full-\(k\) calculations.

\section{Broadening and energy window}

The discrete transition strengths are broadened on an energy grid. If the user
provides an explicit energy window, the grid is constructed from the user
values \(E_{\min}\) and \(E_{\max}\). If no explicit window is provided, the code
determines an automatic transition-energy range from the calculated transitions.
The output reports the actual resolved grid:
\begin{equation}
  \omega_i = E_{\min} + (i-1)\Delta E,\qquad
  i=1,\ldots,N_E,
\end{equation}
with
\begin{equation}
  \Delta E =
  \frac{E_{\max}-E_{\min}}{N_E-1}.
\end{equation}
Here \(N_E\) is the number of energy points requested by the user.

\section{Validation status}

The following implementation checks have been performed:
\begin{itemize}
  \item scalar Fe \(L_3\), \(L_2\), and \(K\) edges give isotropic
  \(x,y,z\) spectra in cubic tests;
  \item two-atom absorber summation was validated for scalar Fe \(L_3\);
  \item first-variation noncollinear SOC full-\(k\) calculations were validated
  for different spin quantization axes;
  \item first-variation noncollinear SOC symmetry-reduced calculations were
  validated against explicit full-\(k\) references for \(L_3\);
  \item a small MPI test reproduced the serial result for the validated
  first-variation noncollinear SOC \(L_3\) case.
\end{itemize}

\section{Relation to DFT+\(U\)}

The XAS matrix-element machinery should also work with DFT+\(U\) in principle.
The reason is that the XAS module uses the eigenvalues, occupations, and
wavefunctions produced by the electronic-structure calculation. DFT+\(U\)
changes the Hamiltonian, and therefore changes \(\varepsilon_{n\kvec}\),
\(f_{n\kvec}\), and the coefficients \(A^{\absr,\sig}_{n\kvec,lm\nu}\), but it
does not require a different dipole matrix-element formula. This point should
still be explicitly validated for representative DFT+\(U\) systems before the
feature is advertised as fully tested.

\section{Physical limitations of the independent-particle approximation}

The present implementation is useful for analyzing orbital character,
polarization dependence, absorber specificity, and trends between calculations.
It is not a complete many-body theory of core-level spectroscopy. The main
limitations are:
\begin{enumerate}
  \item \textbf{No automatic explicit core hole.}
  The formula in Eq.~\eqref{eq:xas-spectrum} uses the electronic structure
  supplied by the calculation. Unless the user performs a separate core-hole
  calculation, the attraction between the excited electron and the core hole is
  not included explicitly.

  \item \textbf{No bound core exciton in the default formula.}
  The excited electron and core hole are not solved as a correlated two-particle
  state.

  \item \textbf{No atomic multiplet physics.}
  Many transition-metal \(L\)-edge spectra contain strong multiplet effects
  arising from local Coulomb interactions between the core hole and valence
  electrons. These are outside the present independent-particle treatment.

  \item \textbf{Broadening is phenomenological.}
  Lifetime and instrumental broadening are represented by the user-selected
  broadening parameter \(\eta\), rather than being computed from first
  principles.

  \item \textbf{Absolute edge energies require alignment.}
  Kohn--Sham eigenvalue differences are not generally quantitative absolute
  excitation energies. Comparison with experiment usually requires an energy
  shift or alignment to a reference.

  \item \textbf{Strong final-state correlations are not included.}
  If the final state is dominated by strong electron-hole attraction, multiplet
  splitting, or other many-body effects, methods beyond this approximation are
  required.
\end{enumerate}

\section{How to read the calculated spectra}

A useful way for students to interpret the calculated spectrum is to ask:
\emph{which unoccupied states are visible from a given core state and photon
polarization?} At the \(K\) edge, a \(1s\) core electron is promoted mainly into
\(p\)-like final states. At the \(L_2\) and \(L_3\) edges, a \(2p\) core electron
is promoted mainly into \(s\)- and \(d\)-like final states. Differences between
the \(x\), \(y\), and \(z\) spectra reveal anisotropy of the unoccupied states
projected onto the absorber atom. In magnetic systems with spin-orbit coupling,
the polarization dependence can also reflect the coupling between orbital
character, spinor structure, and the local magnetic geometry.

The method should therefore be viewed as a transparent independent-particle
probe of orbital character and anisotropy. It is not, by itself, a complete
prediction of experimental core-level line shapes when core-hole attraction,
multiplets, or strong final-state correlations dominate.

\section{Current implementation limitations}

The current prototype has the following deliberate limitations:
\begin{itemize}
  \item only \(K\), \(L_2\), and \(L_3\) are exposed as separate edges;
  \item combined \(L_{23}\) is rejected and should be handled later as a
  separate design step;
  \item the XML interface currently exposes only linear \(x,y,z\) polarizations;
  \item second-variation SOC is guarded and is not part of the validated path;
  \item time-reversal-generated star members for SOC/noncollinear XAS remain
  guarded;
  \item the \texttt{inpgen} noncollinear-SOC \(\theta,\phi\) workflow has a
  known consistency issue: angles used for symmetry generation can be written
  as runtime SOC angles, which FLEUR rejects for noncollinear calculations.
\end{itemize}

\end{document}
